% Copyright (c) 2026 Marcio Luiz Lima de Oliveira and V. el Hakim.
% Distribution and modification are subject to LPPL version 1.3c.
% SPDX-License-Identifier: LPPL-1.3c
% Part of the LPPL-covered work listed in MANIFEST.txt.
% See LICENSE and LICENSING.txt for the licence and scope.

\DeclareRobustCommand{\func}[3]{#1\colon #2\to #3}
% Delimiter helpers retain the existing argument order and default sizing.
% An optional size (e.g. \big or \Big) supports use inside an align row.
\DeclarePairedDelimiter{\ThesisNormDelimiter}{\lVert}{\rVert}
\DeclarePairedDelimiter{\ThesisAbsDelimiter}{\lvert}{\rvert}
\DeclarePairedDelimiterX{\ThesisSetDelimiter}[2]{\lbrace}{\rbrace}{%
  #1\,\delimsize\vert\,#2%
}
\DeclareRobustCommand{\norm}[3][]{%
  \if\relax\detokenize{#1}\relax
    \ThesisNormDelimiter*{#2}%
  \else
    \ThesisNormDelimiter[#1]{#2}%
  \fi
  % An empty index must not add an empty subscript box and its script spacing.
  \if\relax\detokenize{#3}\relax\else_{#3}\fi
}
\DeclareRobustCommand{\Abs}[2][]{%
  \if\relax\detokenize{#1}\relax
    \ThesisAbsDelimiter*{#2}%
  \else
    \ThesisAbsDelimiter[#1]{#2}%
  \fi
}
\DeclareRobustCommand{\Set}[3][]{%
  \if\relax\detokenize{#1}\relax
    \ThesisSetDelimiter{#2}{#3}%
  \else
    \ThesisSetDelimiter[#1]{#2}{#3}%
  \fi
}
\newcommand*{\ctrans}[1]{\xrightarrow{#1}}
\newcommand*{\dtrans}{\xrightarrow{*}}
\DeclareMathOperator{\Lip}{Lip}
\DeclareMathOperator{\Aggr}{Aggr}
\DeclareMathOperator{\supp}{supp}
\DeclareMathOperator*{\argmin}{argmin}
\DeclareMathOperator*{\argmax}{argmax}
\DeclareMathOperator{\diag}{diag}
\DeclareMathOperator{\abs}{\mathbf{abs}}
\DeclareMathOperator{\Skew}{skew}
\DeclareMathOperator{\Prod}{prod}
\DeclareMathOperator{\Vol}{vol}
\DeclareMathOperator{\Aff}{aff}
\DeclareMathOperator{\Conv}{conv}
\DeclareMathOperator{\Int}{int}
\DeclareMathOperator{\Cl}{cl}
\DeclareMathOperator{\Aconv}{dconv}
\DeclareMathOperator{\Tr}{tr}
\DeclareMathOperator{\Span}{span}
\DeclareMathOperator{\Ker}{ker}
\DeclareMathOperator{\Ran}{ran}
\DeclareMathOperator{\Sign}{sign}
\DeclareMathOperator{\Dom}{dom}
\newcommand*{\Orth}[1]{\mathrm{O}(#1)}
\newcommand*{\So}[1]{\mathrm{SO}(#1)}
